%----------------------------------------------------------------------------------------
% БҮЛЭГ 4: ӨГӨГДӨЛ БОЛОВСРУУЛАЛТ
%----------------------------------------------------------------------------------------

\chapter{Өгөгдөл боловсруулалт}

\section{Өгөгдлийн эх сурвалж ба зөвшөөрөл}

Энэхүү судалгаанд хоёр төрлийн өгөгдлийн эх сурвалж ашигласан: нийтийн бенчмарк өгөгдлийн багцууд болон өөрийн цуглуулсан хяналтын камерын тайлбар тайлангийн өгөгдөл.

\subsection{Нийтийн бенчмарк өгөгдлийн багцууд}

\begin{enumerate}
    \item \textbf{RWF-2000 өгөгдлийн сан} \cite{Cheng2021}.
    \newline{\small\sloppy\url{https://github.com/mchengny/RWF2000-Video-Database-for-Violence-Detection}}
    \begin{itemize}
        \item Судалгаанд ашигласан хэмжээ: 1,000 видео клип (500 хүчирхийлэлтэй + 500 хүчирхийлэлгүй).
        \item Эх сурвалж: YouTube-ийн хяналтын камерын бичлэг.
        \item Клипийн урт: Клип тус бүр 5 секунд, 30 fps.
        \item Нягтрал: Хувьсах ($360 \times 288$ -- $1920 \times 1080$).
        \item Лиценз: Судалгааны зориулалтаар чөлөөтэй ашиглах боломжтой.
        \item Онцлог: Бодит хяналтын камерын нөхцөл, олон янзын орчин.
    \end{itemize}

    \item \textbf{Real Life Violence Situations өгөгдлийн багц (зодооны хэсэг)}.
    \newline{\small\url{https://www.kaggle.com/datasets/mohamedmustafa/real-life-violence-situations-dataset}}
    \begin{itemize}
        \item Судалгаанд ашигласан хэмжээ: 500 зодооны клип.
        \item Эх сурвалж: YouTube-ийн бодит хүчирхийллийн бичлэг.
        \item Лиценз: Судалгааны зориулалтаар нийтэд нээлттэй.
        \item Онцлог: Бодит амьдрал дахь хүчирхийллийн олон янзын үзэгдэл.
    \end{itemize}

    \item \textbf{Violence өгөгдлийн багц} ---.
    \begin{itemize}
        \item Судалгаанд ашигласан хэмжээ: 451 хүчирхийлэлгүй видео клип.
        \item Эх сурвалж: Хяналтын камерын болон интернэтийн бичлэг.
        \item Лиценз: Судалгааны зориулалтаар нийтэд нээлттэй.
        \item Онцлог: Олон төрлийн орчин, гэрэлтүүлгийн нөхцөлтэй бичлэг.
    \end{itemize}
\end{enumerate}

\subsection{Өөрийн цуглуулсан өгөгдөл}

\begin{itemize}
    \item \textbf{Эх сурвалж:} Хяналтын камерын тайлбар тайлангийн файлууд.
    \item \textbf{Анхны хэмжээ:} Нийт 2,269 тайлбар тайлан файл цуглуулсан.
    \item \textbf{Агуулга:} Үйл явдлын тайлбар, оролцогчдын дүрслэл, хронологи.
    \item \textbf{Хэл:} Монгол болон англи хосолсон.
\end{itemize}

\subsection{Ёс зүй ба зөвшөөрөл}

Хяналтын камерын системтэй холбоотой өгөгдлийг ашиглахад дараах ёс зүйн зарчмуудыг баримталсан:

\begin{itemize}
    \item Ашигласан бенчмарк өгөгдлийн багцууд нь олон нийтэд нийтлэгдсэн, судалгааны зориулалтаар лицензлэгдсэн.
    \item Хувь хүний таних боломжтой мэдээлэл (нүүр, нэр) боловсруулалтын явцад анонимчлагдсан.
    \item GDPR болон Монгол улсын Хувийн мэдээлэл хамгаалах тухай хуулийн шаардлагуудыг дагаж мөрдсөн \cite{gdpr2016}.
    \item Бичлэгийг зөвхөн зодооны үйл явдал илэрсэн тохиолдолд хадгалах зарчим баримтлав.
\end{itemize}

\section{Өгөгдлийн урьдчилан боловсруулалт}

\subsection{Видео өгөгдлийн боловсруулалт}

Видео бичлэгийг боловсруулахдаа дараах алхмуудыг дагав:

\begin{enumerate}
    \item \textbf{Frame extraction:} Секундэд 1 фрейм буюу оролтын видеоны 30-р фрейм бүрийг сонгосон. Энэ нь тооцооллын нөөцийг хэмнэхийн зэрэгцээ үйл явдлын гол мэдээллийг хадгалах боломжийг олгодог.

    \item \textbf{Фрейм стандартчилал:} Фрейм бүрийг $224 \times 224$ пикселийн хэмжээнд оруулсан.
\end{enumerate}

\subsection{Тайлбар тайлангийн үүсгэлт ба боловсруулалт}

Тайлбар болон тайлан файлууд нь урьдчилан бэлдсэн өгөгдөл биш бөгөөд \textbf{Ollama} орчинд ажиллуулсан \textbf{Qwen3-VL:235b-clouded} загварыг ашиглан хяналтын камерын видео бичлэгийг шууд шинжилж, тухайн шинжилгээний үр дүнд үндэслэн автоматаар үүсгэсэн болно.

\subsubsection*{Тайлбар тайлан үүсгэх үйл явц}

Би Ollama орчинд ажиллаж буй Qwen3-VL:235B загварт зориулан заавар агуулсан Python код боловсруулж, тухайн кодыг ажиллуулан хяналтын камерын видео бичлэгийг оролт болгон өгсөн. Үүний үр дүнд загвар нь видеоны агуулгыг боловсруулж, түүнээс шаардлагатай мэдээллийг нэгтгэн тайлан хэлбэрээр гаргаж өгсөн. Заавар дараах бүтэцтэй байв:

\begin{itemize}
    \item Бичлэгт зодоон болсон эсэхийг тодорхойлох.
    \item Оролцогчдыг дүрслэх (тоо, хувцас, байрлал).
    \item Үйл явдлын хронологийг T0–T3 үе шатаар тайлбарлах.
    \item Нотлох баримт дурдан дүгнэлт, зөвлөмж гаргах.
\end{itemize}

Загвар нь нэг бичлэг бүрт доорх 8 хэсэгтэй бүтэцтэй Монгол хэлний тайлан үүсгэсэн:
\texttt{ШОШГО}, \texttt{ИТГЭЛЦЭЛ}, \texttt{СУБЪЕКТҮҮД}, \texttt{ҮЙЛДЛҮҮД}, \texttt{ЯВЦ}, \texttt{НОТЛОХ БАРИМТ}, \texttt{ДҮГНЭЛТ}, \texttt{ЗӨВЛӨМЖ}.

\subsubsection*{Тайлбар тайлангийн урьдчилан боловсруулалт}

Үүсгэгдсэн тайлангуудыг Q\&A өгөгдлийн багц үүсгэхэд бэлтгэхдээ дараах боловсруулалт хийсэн:

\begin{enumerate}
    \item \textbf{Тэмдэгт цэвэрлэгээ:} Тусгай тэмдэгт, давхардсан зай, HTML тег зэргийг цэвэрлэв.

    \item \textbf{Хэл тодорхойлолт:} Монгол болон англи холимог текстийг тодорхойлж, нэгдмэл боловсруулалт хийв.
\end{enumerate}


\section{Өгөгдлийн багц үүсгэлт}

Багш загвар (Qwen3-VL-30B-A3B) ашиглан нэг тайлбар тайлангаас 4 төрлийн Q\&A хос үүсгэсэн. Анхны 2,269 тайлбар тайлангаас цэвэрлэгээ, шалгалтын дараа сургалтад ашиглах боломжтой тайлангуудыг сонгон авч, нийт ойролцоогоор 9,040 Q\&A хос бүрдүүлсэн:

\begin{enumerate}
    \item \textbf{event\_type} — ``Энэ бичлэгт зодоон болсон уу? ямар эрчимтэй?''.
    \item \textbf{subjects} — ``Үйл явдалд оролцсон хүмүүсийг тайлбарла''.
    \item \textbf{chronology} — ``Үйл явдлын хронологийг тайлбарла (T0--T3)''.
    \item \textbf{conclusion} — ``Нотлох баримт дурдан дүгнэлт гарга''.
\end{enumerate}

\begin{figure}[H]
    \centering
    \includegraphics[width=0.85\textwidth]{\detokenize{Figures/Q&A.png}}
    \caption{Тайлбар тайлангаас Q\&A өгөгдлийн хос үүсгэсэн жишээ}
    \label{fig:qa_generation}
\end{figure}

\section{Train/validation/test хуваалт}

\begin{table}[htbp]
\centering
\caption{Өгөгдлийн хуваалт}
\label{tab:data_split}
\begin{tabular}{lcccc}
\hline
\textbf{Хуваалт} & \textbf{Нийт} & \textbf{Хүчирхийлэл} & \textbf{Хүчирхийлэлгүй} & \textbf{Хувь} \\
\hline
Сургалт (Train)          & 1,560 & 801 (51.3\%)       & 759 (48.7\%)       & 80\% \\
Баталгаажуулалт (Val)    & 195   & 105 (53.8\%)       & 90 (46.2\%)        & 10\% \\
Тест (Test)              & 196   & 94 (48.0\%)        & 102 (52.0\%)       & 10\% \\
\hline
\textbf{Нийт}             & \textbf{1,951} & \textbf{1,000 (51.3\%)} & \textbf{951 (48.7\%)} & \textbf{100\%} \\
\hline
\end{tabular}
\end{table}

Өгөгдлийн багцыг хуваахдаа дараах зарчмуудыг баримталсан:
\begin{itemize}
    \item Нэг видеоноос гаргасан клипүүд нэг өгөгдлийн хуваалтад л байрлах ёстой (өгөгдлийн алдагдлаас сэргийлэх).
    \item Ангилалтын тэнцвэрт байдлыг хадгалах (хүчирхийлэлтэй:хүчирхийлэлгүй харьцааг тогтмол байлгах).
    \item Тест сет нь сургалтын явцад огт ашиглагдаагүй.
\end{itemize}

\section{Өгөгдлийн чанарын үзүүлэлтийн хүснэгт}

\begin{table}[htbp]
\centering
\caption{Өгөгдлийн чанарын үзүүлэлт}
\label{tab:data_quality}
\begin{tabular}{llll}
\hline
\textbf{Үзүүлэлт} & \textbf{Утга} & \textbf{Зорилтот шалгуур} & \textbf{Үнэлгээ} \\
\hline
Нийт өгөгдлийн хэмжээ      & 1,951 клип    & $\geq$ 1,900 & \checkmark \\
Class distribution (хүчирхийлэлтэй) & 51.3\%          & 45--55\%     & \checkmark \\
Class distribution (non-viol)& 48.7\%          & 45--55\%     & \checkmark \\
Missing rate (дутуу утга)    & $<$ 0.1\%     & $<$ 5\%      & \checkmark \\
Duplicate rate (давхардал)   & 0\%           & 0\%          & \checkmark \\
Өгөгдлийн алдагдал (давхцал)       & 0\%           & 0\%          & \checkmark \\
Minimum clip length          & 2 секунд      & $\geq$ 2 с   & \checkmark \\
\hline
\end{tabular}
\end{table}

\section{Өгөгдлийн нэмэгдүүлэлт}

\begin{sloppypar}Сургалтын явцад загварын ерөнхий чадварыг нэмэгдүүлэх зорилгоор дараах augmentation аргуудыг ашигласан:\end{sloppypar}

\begin{itemize}
    \item \textbf{Temporal augmentation:} Клипийн эхлэх цэгийг санамсаргүй байдлаар сонгох.
    \item \textbf{Зааврын хувилбар:} Q\&A зааврын томъёололыг $\pm$3 хувилбараар өөрчлөх.
    \item \textbf{Text augmentation:} Синоним орлуулалт, өгүүлбэрийн дарааллыг өөрчлөх.
\end{itemize}

\begin{sloppypar}Видео фреймүүдэд өнгөний өөрчлөлт, санамсаргүй тайралт болон хэвтээ эргүүлэлт зэрэг стандарт augmentation аргуудыг ашигласан. Харин хяналтын камерын бодит нөхцөлд хэвтээ эргүүлэлт нь орон зайн бодит мэдээллийг гажуудуулах боломжтой тул уг аргыг хязгаарлагдмал хэмжээнд хэрэглэсэн.\end{sloppypar}

%-------------------------------------------------------------------------------
\section{Өгөгдлийн багц ба статистик шинжилгээ}
%-------------------------------------------------------------------------------

\subsection{RWF-2000 өгөгдлийн сангийн дэлгэрэнгүй шинжилгээ}

RWF-2000 (Real World Fight) өгөгдлийн сан нь хяналтын камерын бодит бичлэгээс гаргасан бенчмарк өгөгдөл юм~\cite{Cheng2021}. Энэхүү судалгаанд уг өгөгдлийн сангаас 1,000 клип сонгон ашигласан. Өгөгдлийн сангийн гол шинж чанарууд:

\begin{table}[htbp]
\centering
\caption{RWF-2000 өгөгдлийн сан дэлгэрэнгүй статистик}
\label{tab:rwf2000_stats}
\begin{tabular}{ll}
\hline
\textbf{Үзүүлэлт} & \textbf{Утга} \\
\hline
Судалгаанд ашигласан клип & 1,000 (500 хүчирхийлэлтэй + 500 хүчирхийлэлгүй) \\
Клип бүрийн хугацаа      & 5 секунд (тогтмол) \\
Frame rate                & 30 fps \\
Нийт фрейм               & 150,000 (150 фрейм/клип) \\
Нягтрал (resolution)     & $360 \times 288$ --- $1920 \times 1080$ \\
Дундаж нягтрал            & $640 \times 480$ (ойролцоогоор) \\
Эх сурвалж                & YouTube (нийтийн хяналтын камерын бичлэг) \\
Орчин                    & Гудамж, дотоод орон зай, спортын заал гэх мэт \\
Лиценз                   & CC BY 4.0 (судалгааны зориулалт) \\
Авторын нэр              & Cheng et al. (2021) \\
\hline
\end{tabular}
\end{table}

RWF-2000 нь хяналтын камерын бодит нөхцөлийг сайн дүрсэлдгээрээ онцлогтой: янз бүрийн гэрэлтүүлэг, камерын байрлал, дэвсгэрийн нарийн төвөгтэй байдал зэргийг агуулсан нь загварын ерөнхий чадварыг (generalizability) сайжруулахад тусалдаг.

\subsection{Real Life Violence Situations өгөгдлийн багцын дэлгэрэнгүй шинжилгээ}

Real Life Violence Situations өгөгдлийн багц нь YouTube-ийн бодит хүчирхийллийн бичлэгүүдийг цуглуулан бүтээгдсэн нийтэд нээлттэй өгөгдлийн багц юм. Бид уг өгөгдлийн багцын зодооны хэсгээс 500 клип сонгон ашигласан:

\begin{table}[htbp]
\centering
\caption{Real Life Violence Situations өгөгдлийн багцын статистик}
\label{tab:ucfcrime_stats}
\resizebox{0.95\textwidth}{!}{%
\begin{tabular}{ll}
\hline
\textbf{Үзүүлэлт} & \textbf{Утга} \\
\hline
Судалгаанд ашигласан зодооны клип & 500 \\
Клипийн дундаж хугацаа   & $\approx$4--15 секунд (хувьсах) \\
Бичлэгийн эх сурвалж      & YouTube, нийтийн хяналтын камер \\
Онцлог                    & Бодит амьдрал дахь хүчирхийллийн олон янзын үзэгдэл \\
Лиценз                    & Судалгааны зориулалт (Kaggle нийтлэл) \\
\hline
\end{tabular}
}
\end{table}

Real Life Violence Situations өгөгдлийн багцын клипүүд нь хугацааны хувьд өөр өөр байдаг тул frame extraction алгоритм нь хугацааг уусгасан (normalized) байдлаар 8 фрейм гаргаж авдаг.

\subsection{Violence өгөгдлийн багцын дэлгэрэнгүй шинжилгээ}

Violence өгөгдлийн багцаас тэнцвэржүүлэлтийн зорилгоор 451 хүчирхийлэлгүй клип сонгон ашигласан:

\begin{table}[htbp]
\centering
\caption{Violence өгөгдлийн багцын статистик}
\label{tab:hockey_stats}
\resizebox{0.95\textwidth}{!}{%
\begin{tabular}{ll}
\hline
\textbf{Үзүүлэлт} & \textbf{Утга} \\
\hline
Судалгаанд ашигласан клип & 451 хүчирхийлэлгүй клип \\
Клипийн хугацаа      & $\approx$2--5 секунд \\
Бичлэгийн чанар      & Хувьсах (интернет болон хяналтын камерын бичлэг) \\
Онцлог               & Олон төрлийн орчин, гэрэлтүүлгийн нөхцөл \\
Давуу тал            & Тодорхой тайлбарлагдсан хүчирхийлэлтэй/хүчирхийлэлгүй label \\
\hline
\end{tabular}
}
\end{table}

Violence өгөгдлийн багц нь хүчирхийлэлгүй жишээг нэмэгдүүлж, ангиллын тэнцвэрийг хадгалахад туслах auxiliary өгөгдлийн багц болгон ашигласан.
%-------------------------------------------------------------------------------
\section{Q\&A загварын дизайн ба зорилго}
%-------------------------------------------------------------------------------

\subsection{4 Q\&A промпт сонгосны үндэслэл}

Нэг тайлангаас 4 ялгаатай Q\&A хос үүсгэхэд дараах арга зүйн үндэслэл байна:

\begin{enumerate}
    \item \textbf{Олон талт мэдлэгийн дамжуулалт:} Зөвхөн нэг Q\&A ашиглавал загвар нэг чиглэлийн мэдлэг л эзэмшинэ. Харин 4 ялгаатай чиглэл (ангилалт, субъект, хронологи, дүгнэлт) нь загварт өргөн хүрээний мэдлэг олгодог.
    \item \textbf{Fine-grained supervision:} Зодооны ангилалт, оролцогчдын тодорхойлолт, цаг хугацааны дараалал, дүгнэлтийн хослол нь загварт олон давхаргын дүн шинжилгээ хийх чадварыг суулгадаг.

    \item \textbf{Дүгнэлт гаргалт-ийн уян хатан байдал:} Сургасан загвар нь 4 төрлийн асуултад хариулах чадвартай тул хэрэглэгчийн хэрэгцээнд нийцэх боломжтой.
\end{enumerate}

\subsection{Q\&A загварын жишээ}

Дараах жишээ нь нэг тайлбар тайлангаас үүсгэгдсэн 4 Q\&A хосыг харуулж байна:

\begin{tcolorbox}[tab2, title={Q\&A Жишээ --- event\_type}]
\small
\textbf{Асуулт:} Тайланд дурдагдсан үйл явдалд зодоон, хүчирхийлэл байсан уу? Байгаа бол зодооны зөрчлийн зэрэглэлийг тодорхойлж тайлбарла.\\[4pt]
\textbf{Хариулт:} Тийм, тайланд дунд зэргийн зодооны зөрчил тэмдэглэгдсэн. Хоёр оролцогч хоорондоо зодолдсон бөгөөд нэг нь нөгөөгөө газарт унагаасан байна. Гэмтэл болсон эсэх тодорхойгүй боловч биеийн харилцаа холбоо тогтоогдсон.
\end{tcolorbox}

\begin{tcolorbox}[tab2, title={Q\&A Жишээ --- subjects}]
\small
\textbf{Асуулт:} Үйл явдалд оролцсон хүмүүсийг тодорхойл. Тус бүрийн физик байдал, хувцас, байршлыг нарийвчлан дүрсэл.\\[4pt]
\textbf{Хариулт:} Субъект A: Эрэгтэй хүн, дунд биетэй. Хар өнгийн өмдтэй, цагаан футболктай. Эхний байрлал: зүүн тал. Субъект B: Эрэгтэй хүн, нарийн биетэй. Цэнхэр жинс, улаан дөрвөлжин хээт цамцтай. Эхний байрлал: баруун тал, камераас 3 метрийн зайд.
\end{tcolorbox}

\begin{tcolorbox}[tab2, title={Q\&A Жишээ --- chronology}]
\small
\textbf{Асуулт:} Үйл явдлын цаг хугацааны дарааллыг T0-ээс T3 хүртэл тайлбарла.\\[4pt]
\textbf{Хариулт:} T0 (00:00): Субъект A ба B хоорондоо 2 метрийн зайтай зогсож байна. Маргаан эхэлж байгаа шинж тэмдэг ажиглагдаж байна. T1 (00:02): Субъект B Субъект A руу дөхнө. Хоёрын зай 0.5 м болтол буурна. Субъект A гараа өргөнө. T2 (00:04): Субъект B Субъект A-г мөрөнд нь барин, хажуу руу нь түлхэнэ. T3 (00:06): Субъект A хана руу тулаад тогтворжино. Зодоон намжина.
\end{tcolorbox}

\begin{tcolorbox}[tab2, title={Q\&A Жишээ --- conclusion}]
\small
\textbf{Асуулт:} Нотлох баримтад тулгуурлан зодоон болсон эсэхийг дүгнэж, аюулгүй байдлын зөвлөмж гарга.\\[4pt]
\textbf{Хариулт:} Дүгнэлт: Зодоон болсон. Нотлох баримт: (1) биеийн харилцаа холбоо тогтоогдсон (барих, түлхэх); (2) хоёр субъект хоорондоо дайрсан. Зодооны зөрчлийн зэрэглэл: дунд. Зөвлөмж: Аюулгүй байдлын ажилтанд нэн даруй мэдэгдэж, тухайн орчныг зориудаар шалгана уу. Холбогдох байгаа бол хоёр субъектыг тусгаарлаарай.
\end{tcolorbox}

%-------------------------------------------------------------------------------
\section{Өгөгдлийн хуваарилалтын дэлгэрэнгүй шинжилгээ}
%-------------------------------------------------------------------------------

\subsection{Хуваарилалтын зарчим}

Өгөгдлийн багцыг хуваахдаа дараах нэмэлт зарчмуудыг баримталсан:

\begin{itemize}
    \item \textbf{Видео түвшний хуваарилалт:} Нэг видеоноос авсан бүх клип нэг өгөгдлийн хуваалт дотор байна. Энэ нь нэг видеоны ижил фреймүүд өөр өөр өгөгдлийн хуваалтад орж, алдаатай дүгнэлт гарахаас сэргийлнэ.
    \item \textbf{Давхардал байхгүй байдал:} Test set нь сургалт болон баталгаажуулалтын өгөгдлийн хуваалтад ашиглагдаагүй байсан нь SSIM (Structural Similarity Index) ашиглан хянагдсан.
    \item \textbf{Тэнцвэрт ангилал:} хүчирхийлэлтэй/хүчирхийлэлгүй харьцааг бүх өгөгдлийн хуваалтад ойролцоогоор тэнцүү байлгасан (stratified sampling).
\end{itemize}

\subsection{Хуваарилалтын дэлгэрэнгүй хүснэгт}

\begin{table}[htbp]
\centering
\caption{Өгөгдлийн багц бүрийн өгөгдлийн хуваалтын дэлгэрэнгүй хуваарилалт}
\label{tab:partition_detail}
\resizebox{\textwidth}{!}{%
\begin{tabular}{lccccccc}
\hline
\textbf{Хуваалт} & \textbf{Нийт} & \multicolumn{2}{c}{\textbf{Ангиллын харьцаа}} & \multicolumn{2}{c}{\textbf{Хувь}} & \multicolumn{2}{c}{\textbf{тайлбар}} \\
 & & V & NV & V & NV & V & NV \\
\hline
Сургалт (Train)       & 1,560 & 801 & 759 & 51.3\% & 48.7\% & \multicolumn{2}{c}{80\%} \\
Баталгаажуулалт (Val) & 195   & 105 & 90  & 53.8\% & 46.2\% & \multicolumn{2}{c}{10\%} \\
Тест (Test)           & 196   & 94  & 102 & 48.0\% & 52.0\% & \multicolumn{2}{c}{10\%} \\
\hline
\textbf{Нийт видео} & \textbf{1,951} & \textbf{1,000} & \textbf{951} & \textbf{51.3\%} & \textbf{48.7\%} & \multicolumn{2}{c}{\textbf{100\%}} \\
\hline
\multicolumn{8}{l}{\footnotesize V = хүчирхийлэлтэй, NV = хүчирхийлэлгүй.}
\end{tabular}%
}
\end{table}

\subsection{Нарийвчилсан тохируулга өгөгдлийн багцын хуваарилалт}

Q\&A өгөгдлийн багцын хувьд тайлан түвшний хуваарилалт хийсэн (нэг тайлангийн 4 Q\&A хос нэг өгөгдлийн хуваалт дотор байна):

\begin{table}[htbp]
\centering
\caption{Нарийвчилсан тохируулга Q\&A өгөгдлийн багцын хуваарилалт}
\label{tab:qa_split}
\begin{tabular}{lccc}
\hline
\textbf{Partition} & \textbf{Тайлангийн тоо} & \textbf{Q\&A хос тоо} & \textbf{Хувь} \\
\hline
Сургалт (Train)       & 2,034 & 8,136 & 90\% \\
Баталгаажуулалт (Val) & 226   & 904   & 10\% \\
\textbf{Нийт ашигласан} & \textbf{2,260} & \textbf{9,040} & \textbf{100\%} \\
\hline
\hline
\end{tabular}
\end{table}

Ингэснээр анхны 2,269 тайлбар тайлангаас сургалтын чанарын шалгалтаар тохирсон 2,260 тайланг ашиглаж, нэг тайлангаас 4 Q\&A хос үүсгэснээр нийт 9,040 Q\&A хос бүхий нарийвчилсан тохируулгын өгөгдлийн багц бүрдсэн.

%-------------------------------------------------------------------------------
\section{Өгөгдлийн урьдчилсан боловсруулалтын дэлгэрэнгүй алгоритм}
%-------------------------------------------------------------------------------

\subsection{Фрейм задлах алгоритм}

Видео клипоос 8 фрейм uniform sampling аргаар задалдаг. Секундэд 30 фрейм (30fps) бүхий 5 секундын клипийн хувьд нийт 150 фрейм байдаг. Энэ нь $\lfloor 150 / 8 \rfloor = 18$ фрейм бүрт нэг фрейм задалж, ойролцоогоор 1 секунд бүрт нэг фрейм гаргана:

\begin{equation}
    f_k = \left\lfloor \frac{k \cdot F}{N} \right\rfloor, \qquad k = 0, 1, \ldots, N-1
\end{equation}

Энд $F$ нь видеоны нийт фреймийн тоо, $N = 8$ нь сонгох фреймийн тоо, $f_k$ нь $k$-р фреймийн индекс юм. Энэ арга нь видеоны хугацааны урт өөр байсан ч тогтмол $N=8$ фрейм гаргахыг хангадаг.
